\documentclass[11pt]{article}
\usepackage[a4paper,margin=18mm]{geometry}
\usepackage[no-math]{fontspec}
\setmainfont{Latin Modern Roman}
\usepackage{amsmath,amssymb,amsthm,xfrac,xspace,hyperref,graphicx,comment}
\excludecomment{DefTralics}

\providecommand{\bibliofont}{\small}
\newcommand{\ie}{i.e.,\xspace}
\newcommand{\eg}{e.g.,\xspace}
\newcommand{\etc}{etc.\xspace}
\newcommand{\cf}{cf.\xspace}
\newcommand{\resp}{resp.\xspace}
\newcommand{\smaller}{\small}
\newcommand{\Subsection}{\subsection}
\newcommand{\Subsubsection}{\subsubsection}
\newcommand{\skpt}{\smallskip}
\emergencystretch=3em
\numberwithin{equation}{section}\providecommand{\No}{}
\input{author-preamble.tex}
\makeatletter
\long\def\@makecaption#1#2{\vskip\abovecaptionskip{\small\noindent#1. #2\par}\vskip\belowcaptionskip}
\makeatother
\begin{document}
\begin{center}{\Large Stable item 2039}\end{center}
\paragraph{Source and attribution.}
Naruki Masuda, Hugh Thomas, Andy Tonks, Bruno Vallette, \emph{The diagonal of the associahedra}, Journal de l’École polytechnique — Mathématiques 8 (2021), publisher version, DOI 10.5802/jep.142. \href{https://jep.centre-mersenne.org/articles/10.5802/jep.142/}{Publisher article}. Authors' text: \href{https://creativecommons.org/licenses/by/4.0/}{CC BY 4.0}.
\paragraph{Editorial problem boundary.}
Prove only Theorem 1 below for the standard Loday realisations and the exact transition-map compositions of Definition 12: the resulting cellular nonsymmetric operad and the compatible diagonal maps. All substantive local weighted-facet, oriented-fibre diagonal, continuity and uniqueness constructions are retained. The separate magical-formula Theorem 2 is context only and is not selected.
\paragraph{Difficulty (editorial): H2.}
Face combinatorics of the Loday associahedra does not by itself supply coherent continuous transition maps that commute with the chosen diagonal. The author proves the weighted-facet and oriented-fibre properties needed to construct those maps and establish their continuity and uniqueness across boundary strata. That new global compatibility makes the operadic composition diagrams commute, rather than assuming the final operad axioms from the face labels.
\paragraph{Editorial transcription note.}
Original author language, mathematics and ordering are preserved. Publisher front matter is replaced by this independent article wrapper. Bibliography is recovered from matching cached publisher HTML, including its TeX formulas. No new proof or mathematical repair is supplied. Surrounding original results are retained for conservative context and are not extra selected corpus claims.
\clearpage
\input{author-body.tex}
\begin{thebibliography}{99}
\bibitem{AguiarArdila17} M. Aguiar \& F. Ardila - “Hopf monoids and generalized permutahedra” , 2017
\bibitem{Amorim17} L. Amorim - “The Künneth theorem for the Fukaya algebra of a product of Lagrangians” , Internat. J. Math. 28 (2017) no. 4, article ID 1750026, 38 pages
\bibitem{Brown59} E. H. Brown - “Twisted tensor products. I” , Ann. of Math. (2) 69 (1959), p. 223-246
\bibitem{BilleraSturmfels92} L. J. Billera \& B. Sturmfels - “Fiber polytopes” , Ann. of Math. (2) 135 (1992) no. 3, p. 527-549
\bibitem{BoardmanVogt73} J. M. Boardman \& R. M. Vogt - Homotopy invariant algebraic structures on topological spaces , Lect. Notes in Math. , vol. 347 , Springer-Verlag, Berlin, 1973
\bibitem{CFZ02} F. Chapoton, S. Fomin \& A. Zelevinsky - “Polytopal realizations of generalized associahedra” , Canad. J. Math. 45 (2002) no. 4, p. 537-566
\bibitem{CSZ15} C. Ceballos, F. Santos \& G. M. Ziegler - “Many non-equivalent realizations of the associahedron” , Combinatorica 35 (2015) no. 5, p. 513-551
\bibitem{EilenbergMacLane54} S. Eilenberg \& S. Mac Lane - “On the groups $H(\Pi ,n)$. II. Methods of computation” , Ann. of Math. (2) 60 (1954), p. 49-139
\bibitem{EilenbergZilber53} S. Eilenberg \& J. A. Zilber - “On products of complexes” , Amer. J. Math. 75 (1953), p. 200-204
\bibitem{Forcey08} S. Forcey - “Convex hull realizations of the multiplihedra” , Topology Appl. 156 (2008) no. 2, p. 326-347
\bibitem{GKZ08} I. M. Gelfand, M. M. Kapranov \& A. V. Zelevinsky - Discriminants, resultants and multidimensional determinants , Modern Birkhäuser Classics , Birkhäuser Boston, Inc., Boston, MA, 2008
\bibitem{GaberdielZwiebach97} M. R. Gaberdiel \& B. Zwiebach - “Tensor constructions of open string theories. I. Foundations” , Nuclear Phys. B 505 (1997) no. 3, p. 569-624
\bibitem{Loday04a} J.-L. Loday - “Realization of the Stasheff polytope” , Arch. Math. (Basel) 83 (2004) no. 3, p. 267-278
\bibitem{LodayVallette12} J.-L. Loday \& B. Vallette - Algebraic operads , Grundlehren Math. Wiss. , vol. 346 , Springer-Verlag, Berlin, 2012
\bibitem{May72} J. May - The geometry of iterated loop spaces , Lect. Notes in Math. , vol. 271 , Springer-Verlag, Berlin, 1972
\bibitem{MarklShnider06} M. Markl \& S. Shnider - “Associahedra, cellular $W$-construction and products of $A_\infty $-algebras” , Trans. Amer. Math. Soc. 358 (2006) no. 6, p. 2353-2372
\bibitem{MSS} M. Markl, S. Shnider \& J. D. Stasheff - Operads in algebra, topology and physics , Math. Surveys and Monographs , vol. 96 , American Mathematical Society, Providence, RI, 2002
\bibitem{Proute11} A. Prouté - “$A_\infty $-structures. Modèles minimaux de Baues-Lemaire et Kadeishvili et homologie des fibrations” , Repr. Theory Appl. Categ. (2011) no. 21, p. 1-99, Reprint of the 1986 original, With a preface to the reprint by Jean-Louis Loday
\bibitem{Seidel08} P. Seidel - Fukaya categories and Picard-Lefschetz theory , Zürich Lectures in Advanced Math. , European Mathematical Society, Zürich, 2008
\bibitem{Serre51} J.-P. Serre - “Homologie singulière des espaces fibrés. Applications” , Ann. of Math. (2) 54 (1951), p. 425-505
\bibitem{StasheffShnider97} J. D. Stasheff \& S. Shnider - “From operads to “physically” inspired theories (Appendix B)” , in Operads: Proceedings of Renaissance Conferences (Hartford, CT/Luminy, 1995) , Contemp. Math. , vol. 202 , American Mathematical Society, Providence, RI, 1997, p. 53-81
\bibitem{Stasheff63} J. D. Stasheff - “Homotopy associativity of $H$-spaces. I \& II” , Trans. Amer. Math. Soc. 108 (1963), p. 275-292 \& 293–312
\bibitem{Stasheff70} J. D. Stasheff - $H$-spaces from a homotopy point of view , Lect. Notes in Math. , vol. 161 , Springer-Verlag, Berlin, 1970
\bibitem{Steenrod47} N. E. Steenrod - “Products of cocycles and extensions of mappings” , Ann. of Math. (2) 48 (1947), p. 290-320
\bibitem{SaneblidzeUmble04} S. Saneblidze \& R. Umble - “Diagonals on the permutahedra, multiplihedra and associahedra” , Homology Homotopy Appl. 6 (2004) no. 1, p. 363-411
\bibitem{Tamari51} D. Tamari - Monoïdes préordonnés et chaînes de Malcev , Thèse de Mathématique, Paris, 1951
\bibitem{Ziegler95} G. M. Ziegler - Lectures on polytopes , Graduate Texts in Math. , vol. 152 , Springer-Verlag, New York, 1995
\end{thebibliography}
\end{document}
